\documentclass[noanswers, addpoints, 11pt]{exam}
\usepackage[utf8]{inputenc}
\usepackage[french]{babel}
\usepackage[T1]{fontenc}
\usepackage{xcolor}
\usepackage{listings}
\usepackage{graphicx}
\usepackage[most]{tcolorbox}
\usepackage{multicol} % Nécessaire pour les deux colonnes

% Définition des couleurs comme dans l'exemple
\definecolor{codegreen}{rgb}{0,0.6,0}
\definecolor{codegray}{rgb}{0.5,0.5,0.5}
\definecolor{codepurple}{rgb}{0.58,0,0.82}
\definecolor{backcolour}{rgb}{0.95,0.95,0.92}
 
\lstdefinestyle{mystyle}{
    backgroundcolor=\color{backcolour},   
    commentstyle=\color{codegreen},
    keywordstyle=\color{magenta},
    numberstyle=\tiny\color{codegray},
    stringstyle=\color{codepurple},
    basicstyle=\ttfamily\footnotesize, % Police monospace pour le code
    breakatwhitespace=false,         
    breaklines=true,                 
    captionpos=b,                    
    keepspaces=true,        
    numbers=none, % Pas de numéro de ligne pour ces courts snippets                   
    numbersep=5pt,                  
    showspaces=false,                
    showstringspaces=false,
    showtabs=false,                 
    tabsize=4
}
\lstset{style=mystyle, language=Python}

% Informations du contrôle
\title{Contrôle de rattrapage : Algorithmique en Python}
\author{Babash - NSI Tani Malandi et Johanne Bakalara}
\date{Novembre 2025, Durée: 1h}

\begin{document}
\runningfooter{}{\footnotesize Babash - NSI Tani Malandi et Johanne Bakalara \textperiodcentered{} CC BY-SA 4.0 \textperiodcentered{} page \thepage}{}
\firstpagefooter{}{\footnotesize Babash - NSI Tani Malandi et Johanne Bakalara \textperiodcentered{} CC BY-SA 4.0 \textperiodcentered{} page \thepage}{}

\maketitle

% Configuration En-tête et Pied de page
\pagestyle{headandfoot}
\thispagestyle{headandfoot}
\firstpageheader{Nom : \makebox[4cm]{\dotfill} \\  Prénom : \makebox[4cm]{\dotfill}\\ Classe : \makebox[4cm]{\dotfill}}{}{}

% Consignes (Style boîte encadrée)
\section*{Consignes:}
\begin{center} 
\fbox{\parbox{0.9\linewidth}{
\begin{center}
\textbf{Interdiction d'utiliser la calculatrice, le téléphone ou l'ordinateur.}\\
Pour chaque question, analysez le code Python et écrivez dans la case réponse ce que le programme affiche \textbf{à la fin de l'exécution}.\\
Si le programme ne s'arrête jamais (boucle infinie) ou provoque une erreur, écrivez-le explicitement.
\end{center}
}}
\end{center}

% Tableau de notation
\section*{Barème :}
\begin{center}
  \par\bigskip
  \gradetable[h][questions]
\end{center}

\section*{Questions}

\begin{questions}


\question Compréhension élémentaire des structures de contrôle
\begin{multicols}{2}
\begin{parts}
% Q1
\part[1] Variable et réaffectation
\begin{lstlisting}
x = 10
x = x + 5
x = 2
print(x)
\end{lstlisting}
\textbf{Réponse :} \dotfill \vfill

% Q2
\part[1] Condition simple
\begin{lstlisting}
age = 15
if age > 18:
    print(1)
else:
    print(0)
\end{lstlisting}
\textbf{Réponse :} \dotfill \vfill

% Q3
\part[1] Structure Elif
\begin{lstlisting}
n = 50
if n < 20:
    res = 1
elif n < 60:
    res = 2
else:
    res = 3
print(res)
\end{lstlisting}
\textbf{Réponse :} \dotfill \vfill

% Q4
\part[1] Ifs consécutifs
\begin{lstlisting}
valeur = 10
if valeur > 5:
    valeur = valeur - 5
if valeur == 5:
    print(100)
else:
    print(0)
\end{lstlisting}
\textbf{Réponse :} \dotfill \vfill

% Q5
\part[1] Boucle For simple
\begin{lstlisting}
s = 0
for k in range(3):
    s = s + k
print(s)
\end{lstlisting}
\textbf{Réponse :} \dotfill \vfill

% Q6
\part[1] Boucle While simple
\begin{lstlisting}
compteur = 3
while compteur > 0:
    compteur = compteur - 1
print(compteur)
\end{lstlisting}
\textbf{Réponse :} \dotfill \vfill

% Q7
\part[1] Analyse de boucle (Piège)
\begin{lstlisting}
a = 10
while a > 5:
    a = a + 1
print(a)
\end{lstlisting}
\textbf{Réponse :} \dotfill \vfill

\end{parts}
\end{multicols}

\question Opérateurs logiques
\begin{multicols}{2}
\begin{parts}
% Q8
\part[1] Opérateur logique AND
\begin{lstlisting}
x = 10
y = 5
if x > 5 and y < 2:
    print("A")
else:
    print("B")
\end{lstlisting}
\textbf{Réponse :} \dotfill \vfill

% Q9
\part[1] Opérateurs OR et NOT
\begin{lstlisting}
test = False
if not test or 5 == 5:
    print("Oui")
else:
    print("Non")
\end{lstlisting}
\textbf{Réponse :} \dotfill \vfill

\end{parts}
\end{multicols}

\question Fonction
\begin{multicols}{2}
\begin{parts}
% Q10
\part[1] Définition de fonction
\begin{lstlisting}
def calcul(x):
    return x * 3

res = calcul(4)
print(res)
\end{lstlisting}
\textbf{Réponse :} \dotfill \vfill

% Q11
\part[1] Appels imbriqués
\begin{lstlisting}
def f1(n):
    return n + 1

def f2(n):
    return f1(n) * 2

print(f2(3))
\end{lstlisting}
\textbf{Réponse :} \dotfill \vfill

\part[1] Ordre des affichages
\begin{lstlisting}
def fct_a():
    print("A")

def fct_b():
    print("B")
    fct_a()

fct_b()
\end{lstlisting}
\textbf{Réponse :} \dotfill \vfill

\end{parts}
\end{multicols}

\newpage
\question Questions de synthèse 
\begin{multicols}{2}
\begin{parts}
% Q13
\part[1] Booléens et If successifs
\begin{lstlisting}
actif = True
score = 0
if actif:
    score = 10
    actif = False
if actif:
    score = 20
print(score)
\end{lstlisting}
\textbf{Réponse :} \dotfill \vfill

\vspace{1cm}
% Q14
\part[1] Range avec bornes
\begin{lstlisting}
acc = 0
for i in range(2, 5):
    acc = acc + i
print(acc)
\end{lstlisting}
\textbf{Réponse :} \dotfill \vfill

\end{parts}
\end{multicols}

\question Questions de synthèse avancé
\begin{multicols}{2}
\begin{parts}
% Q12
\part[1] Logique imbriquée
\begin{lstlisting}
i = 0
c = 0
while i < 3:
    if i < 2:
        c = c + 10
    i = i + 1
print(c)
\end{lstlisting}
\textbf{Réponse :} \dotfill \vfill
% Q15
\part[1] Synthèse boucle conditionnelle
\begin{lstlisting}
total = 0
for i in range(5):
    if i > 2:
        total = total + 1
    else:
        total = total + 2
print(total)
\end{lstlisting}
\textbf{Réponse :} \dotfill \vfill

% Q16
\part[1] Synthèse : Évolution variable
\begin{lstlisting}
a = 0
for i in range(5):
    if i < 3:
        a = a + 2
    else:
        a = a - 1
print(a)
\end{lstlisting}
\textbf{Réponse :} \dotfill \vfill

% Q17
\part[1] Synthèse : Fonction booléenne
\begin{lstlisting}
def est_grand(x):
    if x > 10:
        return True
    else:
        return False

n = 5
if est_grand(n + 10):
    print(1)
else:
    print(0)
\end{lstlisting}
\textbf{Réponse :} \dotfill \vfill

% Q18
\part[1] Synthèse : While complexe
\begin{lstlisting}
n = 6
res = 0
while n > 0:
    if n > 3:
        res = res + n
    n = n - 2
print(res)
\end{lstlisting}
\textbf{Réponse :} \dotfill \vfill

% Q19
\part[1] Synthèse : Fonction dans boucle
\begin{lstlisting}
def transform(x):
    if x > 5:
        return 0
    return x

total = 0
for k in range(4, 7):
    total = total + transform(k)
print(total)
\end{lstlisting}
\textbf{Réponse :} \dotfill \vfill

\end{parts}
\end{multicols}
\end{questions}

\end{document}